% Source: https://artofproblemsolving.com/community/c3h1870135_question_about_intro_to_alg_texts825
% Stable identity: aops-topic-1870135-q01
Suppose, for contradiction, that a point $(x_0,y_0)$ lies on both lines:
\[
y_0-y_1=m(x_0-x_1)\qquad\text{and}\qquad y_0-y_2=m(x_0-x_2).
\]
Subtracting the second equation from the first gives
\[
y_2-y_1=m(x_2-x_1).
\]
This says that $P_2$ satisfies the equation of $\ell_1$.

Now take any point $(x,y)$ of $\ell_2$, so $y-y_2=m(x-x_2)$. Adding $y_2-y_1=m(x_2-x_1)$ gives $y-y_1=m(x-x_1)$, so $(x,y)$ lies on $\ell_1$. Hence $\ell_2\subseteq\ell_1$. The same argument with the roles exchanged gives $\ell_1\subseteq\ell_2$, so $\ell_1=\ell_2$. This contradicts the assumption that the lines are distinct.

\begin{asy}
usepackage("amsmath");
size(6cm);
defaultpen(fontsize(10pt));
import graph;
draw((-0.5,0)--(5.5,0),gray,Arrow(5));draw((0,-0.5)--(0,4.6),gray,Arrow(5));
real m=0.5;
pair P1=(2.6,2.8),P2=(3,0.5);
draw((-0.4,P1.y+m*(-0.4-P1.x))--(5.3,P1.y+m*(5.3-P1.x)),blue+linewidth(1));
draw((-0.4,P2.y+m*(-0.4-P2.x))--(5.3,P2.y+m*(5.3-P2.x)),red+linewidth(1));
dot(P1);dot(P2);
label("$P_1(x_1,y_1)$",P1,NW);label("$P_2(x_2,y_2)$",P2,SE);
label("$\ell_1$",(5.3,P1.y+m*(5.3-P1.x)),E,blue);label("$\ell_2$",(5.3,P2.y+m*(5.3-P2.x)),E,red);
label("$x$",(5.5,0),S);label("$y$",(0,4.6),W);
label("\mbox{}",(min(currentpicture,user=true)+max(currentpicture,user=true))/2,black); // keeps the text after the figure black when LaTeX typesets the labels inline
\end{asy}

Therefore two distinct lines with the same slope have no common point: they are parallel.
